\section{Three receivers}\label{sec:receivers}

A pattern is only as credible as its known uses. This section presents three: one application,
\emph{taskboard}, implemented three times with identical behaviour. We call them
\emph{receivers} because they do not define the Assertion Ladder; they take it on, each in the
idiom of its language, and each gains something from it. The paper works backwards from those
gains to the capability they share.

\subsection{The application}

Taskboard is a Kanban API with one entity. A task moves from \emph{todo} to \emph{in progress}
to \emph{done}, and may move back from \emph{in progress} to \emph{todo}. No more than a
configured number of tasks may be in progress at once. When a task is completed the team is
told. Every refused request is answered with an RFC~9457 problem document carrying a stable
error code. The application is deliberately small: 283 lines of application code in Java, 223
in Python and 242 in TypeScript, down from 637, 729 and 957 in a first iteration that used
dependency-injection containers, object-relational mappers and a second entity
\citep{griffin2026convention}. Small enough to read whole is a property the ladder depends on:
the \pattern{Worked Slice} only works if the slice can be read in one sitting.

Each receiver is organised in the same six layers: \repo{api}, \repo{services},
\repo{repositories}, \repo{domain}, and the side layers \repo{events} and \repo{config}, with one
composition root. Imports point only downwards, the domain imports nothing but the language's
standard library, and each outside capability (HTTP, SQL, the environment, logging) belongs to
exactly one layer.

\subsection{A steel thread through all three kinds of intent}

Figure~\ref{fig:thread} follows one rule, the work-in-progress limit, from its requirement to
the gate. The requirement is written once in EARS:

\begin{quote}\small
\textbf{R-WIP-1} (state-driven). While \texttt{WIP\_LIMIT} tasks are in progress, when another
task is moved to \texttt{in\_progress}, the board shall refuse with
\texttt{wip\_limit\_exceeded} and leave the task where it was.
\end{quote}

\noindent Its pattern predicts a service method, because the rule needs to count other tasks.
The scenario that covers it cites its id and reads as given / when / then
(Listing~\ref{lst:scenario}). The domain error it raises carries a code and no HTTP; one handler
maps it to a status, and a closure check refuses a new error without one. The limit arrives
through a constructor from the one settings object. Every step is checked by the same command.

\begin{figure}[t]
  \centering
  % One rule, the work-in-progress limit, through every kind of intent and every layer.
\begin{tikzpicture}[
    font=\footnotesize,
    step/.style={draw=#1, line width=0.9pt, fill=#1!7, rounded corners=2pt, text width=3.05cm,
                 minimum height=1.35cm, align=center, inner sep=4pt},
    arr/.style={-{Stealth[length=2.3mm]}, line width=0.8pt, muted},
  ]
  \node[step=requirement] (req) {\textbf{Requirement}\\R-WIP-1, state-driven\\\texttt{REQUIREMENTS.md}};
  \node[step=behaviour, right=0.45cm of req] (sc) {\textbf{Scenario}\\given a full board\ldots\\cites R-WIP-1};
  \node[step=structure, right=0.45cm of sc] (svc) {\textbf{Service}\\guard in \texttt{move\_task}\\(needs other tasks)};
  \node[step=structure, right=0.45cm of svc] (err) {\textbf{Domain error}\\\texttt{WipLimitExceeded}\\code, no HTTP};
  \node[step=structure, below=0.9cm of err] (map) {\textbf{One translation}\\error $\to$ 409\\exhaustive map};
  \node[step=structure, left=0.45cm of map] (cfg) {\textbf{Setting}\\\texttt{WIP\_LIMIT}\\read once, injected};
  \node[step=muted, left=0.45cm of cfg, text width=6.55cm] (gate) {\textbf{Single gate}: compiler, architecture rules,\\requirement check, scenarios, contract, API tests};
  \draw[arr] (req) -- (sc);
  \draw[arr] (sc) -- (svc);
  \draw[arr] (svc) -- (err);
  \draw[arr] (err) -- (map);
  \draw[arr] (svc.south) |- ($(cfg.north)+(0,0.3)$) -- (cfg.north);
  \draw[arr] (map) -- (cfg);
  \draw[arr] (cfg) -- (gate);
\end{tikzpicture}

  \caption{The work-in-progress limit as a steel thread. Colours mark the kind of intent each
  step carries; every step is checked by the single gate.}
  \label{fig:thread}
\end{figure}

\begin{lstlisting}[language=Python,caption={A scenario in the Python receiver. The docstring cites the requirement and reads given / when / then; the requirement check fails if either is missing.},label=lst:scenario]
def test_the_wip_limit_is_enforced(service: TaskService) -> None:
    """R-WIP-1: given a full board, when another task is started,
    then it is refused and stays in todo."""
    first, second = service.create_task("First"), service.create_task("Second")
    service.move_task(first.id, Status.IN_PROGRESS)
    with pytest.raises(WipLimitExceeded):
        service.move_task(second.id, Status.IN_PROGRESS)
    assert service.get_task(second.id).status is Status.TODO
\end{lstlisting}

\subsection{Binding paths}

Table~\ref{tab:binding} shows how each receiver binds every rung pattern. Reading a row shows
one pattern bound three ways; reading a column shows one receiver's complete path.

\begin{table}[t]
  \centering\small
  \caption{Binding paths of the three receivers.}
  \label{tab:binding}
  \begin{tabularx}{\linewidth}{@{}>{\raggedright\arraybackslash}p{3.1cm} X X X@{}}
    \toprule
    Rung pattern & Java (Spring Boot) & Python (FastAPI) & TypeScript (Fastify) \\
    \midrule
    \pattern{Worked Slice} & \repo{moveTask}, service to controller & \repo{move\_task}, service to route & \repo{moveTask}, service to route \\
    \pattern{Short Manifest} & \repo{AGENTS.md}, 480 words & \repo{AGENTS.md}, 427 words & \repo{AGENTS.md}, 468 words \\
    \pattern{Executable Rule} & ArchUnit over bytecode, 8 rules & AST walk, 7 rules incl.\ strict type check & source text, 5 rules \\
    \pattern{Closure Check} & errors by the compiler; requirements by \repo{RequirementsTest} & errors and requirements by tests & errors by the compiler; requirements by a test \\
    \pattern{Honest Double} & contract run on JDBC and in-memory & contract run on SQLite and in-memory & contract run on SQLite and in-memory \\
    \pattern{Teaching Failure} & \repo{because(\ldots)} on every rule & a message on every assertion & a message on every assertion \\
    \pattern{Unrepresentable Mistake} & sealed \repo{DomainException}, exhaustive \repo{switch} & not available; bound at L3 & \repo{Record<ErrorCode, number>} \\
    \pattern{Single Gate} & \repo{mvn verify} & \repo{pytest} & \repo{npm test} \\
    \pattern{Human Boundary} & \multicolumn{3}{>{\raggedright\arraybackslash}p{\dimexpr\linewidth-3.8cm\relax}}{the same three items in each manifest: dependencies; the rules and manifest; schema and endpoint shape} \\
    \bottomrule
  \end{tabularx}
\end{table}

\paragraph{Java} The Java receiver reaches the top of the ladder most often. Its domain errors
form a sealed hierarchy, and the handler that maps them to HTTP uses a \texttt{switch} with no
default branch, so a new error does not compile until it has a status. Its architecture rules
are ArchUnit rules over bytecode, the most precise of the three enforcement styles, and each
carries a \texttt{because} clause that ArchUnit prints on failure (Listing~\ref{lst:archunit}).
Its requirement check uses ArchUnit to import the test classes themselves and asserts that every
scenario's \texttt{@DisplayName} opens with a requirement id.

\begin{lstlisting}[language=Java,caption={An executable rule with a teaching failure in the Java receiver. It closes a gap that no behavioural test can see (Section~\ref{sec:evaluation}).},label=lst:archunit]
@ArchTest
static final ArchRule onlyEventHandlersHaveSideEffects = noClasses()
    .that().resideOutsideOfPackage("..events..")
    .should().dependOnClassesThat().resideInAnyPackage(
        "org.slf4j..", "java.util.logging..", "org.apache.commons.logging..")
    .because("side effects (logs, mail, audit) listen for a domain "
           + "event; publish one instead");
\end{lstlisting}

\paragraph{Python} Python cannot make an unmapped error unrepresentable, so the Python receiver
binds that rule one rung lower: a \pattern{Closure Check} asserts that every subclass of
\texttt{DomainError} has an entry in the status map. It compensates elsewhere. Its architecture
rules walk the syntax tree, which lets them express rules the other receivers cannot, such as
``a route contains no \texttt{if}, loop, comprehension, \texttt{raise} or \texttt{try}''. And the
strict type checker runs inside the same gate, so type errors are reported as failures of the
single command rather than of a separate tool.

\paragraph{TypeScript} The TypeScript receiver reaches L5 through the type system: the status
map is a \texttt{Record} keyed by a closed union of error codes, so a missing entry is a
compile error (Listing~\ref{lst:record}). Its architecture rules are the cheapest of the three,
regular expressions over source text, and the least precise; they are adequate for a codebase
this small and would need a parser for a larger one.

\begin{lstlisting}[caption={Unrepresentable Mistake in the TypeScript receiver.},label=lst:record]
// Record<ErrorCode, number> is exhaustive: a new error code
// will not compile until it has a status.
const STATUS_BY_CODE: Record<ErrorCode, number> = {
  task_not_found: 404,
  invalid_transition: 409,
  wip_limit_exceeded: 409,
};
\end{lstlisting}

\subsection{Layered, not Clean, and what that costs}

All three receivers are layered rather than hexagonal or Clean. The persistence interface lives
beside its implementation in the repositories layer rather than as a port owned by the core, so
dependency inversion does not stop a service from reaching SQL. Capability confinement does:
a rule that only the repositories layer may use the database driver, only the api layer the web
framework, only the config layer the environment, and only the events layer logging. We read
this as a finding about the ladder rather than about architecture styles. A handful of
\pattern{Executable Rule}s recovers most of the domain protection Clean Architecture provides,
without the ports, adapters and interactors an agent would otherwise have to learn and
reproduce in every feature. The trade is ceremony against rules; the ladder lets a team make it
explicitly.

\subsection{What each receiver gains}

Working backwards from the receivers, each gains the same four things. The domain stays in one
place, which the seeding results in Section~\ref{sec:evaluation} measure directly. A new feature
has a recipe: state the rule in EARS, let its pattern choose the layer, copy the slice, run the
gate. Failures say what to do. And the manifest a new contributor reads is under 500 words,
because everything checkable lives in the checks. The receivers differ only in which rung
carries which rule, and those differences are themselves instructive: the same rule set,
bound by different languages, lands on different rungs.
